\documentclass[10pt,oneside,openany]{memoir}
\input{/Users/gcrescenzo/Documents/School/LaTeX Documents/latex-class-notes/preamble-hw.tex}
\input{/Users/gcrescenzo/Documents/School/LaTeX Documents/latex-class-notes/makros.tex}
\begin{document}
\homeworktitle
{Math 5023} %class
{Homework 2} %assignment
{Gianluca Crescenzo} %name

%%%%%%%%%%%%%%%%%%%%%%%%%%%%%%%%%%%
%%%%%%%%%%%%%%%%%%%%%%%%%%%%%%%%%%%
%%%%%%%%%%%%%%%%%%%%%%%%%%%%%%%%%%%

{\color{red} I skipped problem 8 because I had Covid this week and did not manage my time properly enough.}



\begin{exercise}
For an $R$-module $M$, let
$\operatorname{End}_R(M)=\operatorname{Hom}_R(M,M)$, the $R$-module of
\textit{endomorphisms}. We have seen that it is a ring. Compute
$\operatorname{End}_R(M)$ when $M=R$ and $M=R/I$ for an ideal $I$.
\end{exercise}
	{\color{blue} \begin{proof}
Claim: $\End_R(R/I) \cong R/I$. Define $\varphi:R \rightarrow \End_R(R/I)$ by $r \mapsto L_{r+I}$, where $L_{r+I}:R/I \rightarrow R/I$ is left multiplication by $r+I$. For $r_1,r_2 \in R$ we have:
\begin{align*}
	\varphi(r_1 + r_2)(x+I)
	&= L_{r_1+r_2+I}(x+I) \\
	&= (r_1+r_2)x+I \\
	&= (r_1x+I)+(r_2x+I) \\
	&= \varphi(r_1)(x+I)+\varphi(r_2)(x+I) \\
	&= \bigl(\varphi(r_1)+\varphi(r_2)\bigr)(x+I).
\end{align*}
For $r \in R$ and $a \in R$ we have:
\begin{align*}
	\varphi(ar)(x+I)
	&= L_{ar+I}(x+I) \\
	&= arx+I \\
	&= a(rx+I) \\
	&= a\varphi(r)(x+I).
\end{align*}
Thus $\varphi$ is an $R$-module homomorphism. By the First Isomorphism Theorem for Modules, $R/\ker\varphi \cong \varphi(R)$. For any $f \in \End_R(R/I)$, note that
\[
f(x+I)=f(1+I)\cdot x=L_{f(1+I)}(x+I),
\]
Hence $\varphi\bigl(f(1+I)\bigr)=f$. Similarly:
\begin{align*}
	\ker\varphi
	&= \{r\in R:\varphi(r)(x+I)=0+I \text{ for all }x\in R\} \\
	&= \{r\in R:L_{r+I}(x+I)=0+I \text{ for all }x\in R\} \\
	&= \{r\in R:rx\in I \text{ for all }x\in R\} \\
	&= I.
\end{align*}
Thus $R/I\cong\End_R(R/I)$. Taking $I = \{0\}$, we also have $R \cong \End_R(R)$.
	\end{proof}}

\begin{exercise}
We say an $R$-module has \textit{rank one} if it has a linearly independent
generating set consisting of a single element. Any ring $R$ with $1$ can be
viewed as an $R$-module, namely the $R$-module $\langle 1_R\rangle$. We say
it is a ``module over itself''. Consider $R=\mathbb{R}[x,y]$, viewed as a
module over itself. Prove that it is free of rank one, but that the submodule
$I=(x,y)$ is not free, and needs at least two generators.
\end{exercise}
	{\color{blue} \begin{proof}
			Since $\left<1_{R[x,y]} \right> = R[x,y]$, $R[x,y]$ is rank one with basis $\{1_{R[x,y]}\}$. This is indeed a basis, as it cannot be linearly dependent. Hence it is free of rank one. Moreover, the submodule $I = (x,y)$ is not principal, hence it cannot be generated by a single element. It remains to show that $I$ is not free. Note that $\{x,y\}$ is not linearly independent as $(-y)x + (x)y = 0$. Thus $I$ is not free. 
	\end{proof}}

\begin{exercise}
Let $I\subset R$ be an ideal. Prove that $I$ is a free $R$-module if and only
if $I$ is principal that is generated by a non zero-divisor.
\end{exercise}
	{\color{blue}\begin{proof}
		($\Rightarrow$) Let $I$ be a free module. Claim: $I$ is rank 1. Indeed, if $I$ had $e_1$ and $e_2$ as basis elements, then $e_2e_1 - e_1e_2 = 0$ by the commutativity of $R$. This means $e_1$ and $e_2$ are linearly dependent, a contradiction. Since $I$ is rank 1, $I = (e)$ for some $e \in R$. Moreover, $e$ is not a zero divisor, because if it were then there would exist and element $r \in R$, $r \neq 0$ such that $re = 0$. This means $e$ is linearly dependent, again contradicting the fact that $I$ is free. Thus $I$ is a principal ideal generated by a non zero-divisor.

		($\Leftarrow$) Suppose $I = (e)$ where $e$ is not a zero divisor. Then $e$ is linearly independent, hence $I$ is free of rank $1$.
	\end{proof}}

\begin{exercise}
An \emph{idempotent} $e$ of a ring $R$ (with $1$) is any element satisfying
$e^2=e$. We usually assume an idempotent is nontrivial, meaning $e\ne 0,1$.
Suppose $R$ has a nontrivial idempotent $e$. Prove that there is an $R$-module
isomorphism $Re\oplus R(1-e)\simeq R$.
\end{exercise}
	{\color{blue} \begin{proof}
		Define $\varphi:R \rightarrow Re \oplus R(1-e)$ by $r \mapsto re + r(1-e)$. For $r_1,r_2 \in R$ we have:
		\begin{align*}
			\varphi(r_1 + r_2)
			&= (r_1 + r_2)e + (r_1 + r_2)(1-e) \\
			&= r_1e + r_1(1-e) + r_2e + r_2(1-e) \\
			&= \varphi(r_1) + \varphi(r_2).
		\end{align*}
		Similarly, for $r,a \in R$:
		\begin{align*}
			\varphi(ar)
			&= are + ar(1-e) \\
			&= a(re + r(1-e)) \\
			&= a\varphi(r).
		\end{align*}
		Thus $\varphi$ is an $R$-module homomorphism. For $r_1,r_2 \in R$, suppose $\varphi(r_1) = \varphi(r_2)$. Then $r_1e + r_1(1-e) = r_2e + r_2(1-e)$, which is equivalent to $r_1=r_2$. Thus $\varphi$ is injective. Suppose now that $x \in Re \oplus R(1-e)$. Then $x=a+b$, where $a \in Re$ and $b \in R(1-e)$. Then $\varphi(a+b) = \varphi(a) + \varphi(b) = ae + a(1-e) + be + b(1-e) = a+b = x$. Thus $\varphi$ is surjective, which means $R \cong Re \oplus R(1-e)$.

	\end{proof}}

\begin{exercise}
Let $A,B,M,N$ be $R$-modules. Use the universal properties of the coproduct
and product to prove:
\begin{enumerate}
\item[(a)] There is a natural isomorphism
$\operatorname{Hom}_R(A\oplus B,M)\longrightarrow
\operatorname{Hom}_R(A,M)\times\operatorname{Hom}_R(B,M)$.

\item[(b)] There is a natural isomorphism
$\operatorname{Hom}_R(A,M\times N)\longrightarrow
\operatorname{Hom}_R(A,M)\times\operatorname{Hom}_R(A,N)$.
\end{enumerate}
\end{exercise}
	{\color{blue} \begin{proof}
		(a) Since coproducts exist in the category of $R$-modules, let $\iota_A:A \rightarrow A \oplus B$ and $\iota_B : B \rightarrow A \oplus B$ be the natural inclusions. Define $\varphi:\Hom_R(A \oplus B,M) \rightarrow \Hom_R(A,M) \times \Hom_R(B,M)$ by $\varphi(\tau) := (\tau \circ \iota_A, \tau \circ \iota_B)$. For $\tau_1,\tau_2 \in \Hom_R(A\oplus B,M)$, we have:
		\begin{align*}
			\varphi(\tau_1 + \tau_2)
			&= \bigl( (\tau_1 + \tau_2)\circ \iota_A, (\tau_1 + \tau_2)\circ \iota_B \bigr) \\
			&= (\tau_1 \circ \iota_A, \tau_1 \circ \iota_B) + (\tau_2 \circ \iota_A, \tau_2 \circ \iota_B) \\
			&= \varphi(\tau_1) + \varphi(\tau_2).
		\end{align*}
		For $\tau \in \Hom_R(A\oplus B,M)$ and $r \in R$ we have:
		\begin{align*}
			\varphi(r\tau)
			&= \bigl( (r\tau)\circ \iota_A, (r\tau) \circ \iota_B \bigr) \\
			&= r(\tau \circ \iota_A, \tau \circ \iota_B) \\
			&= r\varphi(\tau).
		\end{align*}
		Thus $\varphi$ is an $R$-module homomorphism. It remains to show that $\varphi$ is injective and surjective. For $\tau_1,\tau_2 \in \Hom_R(A \oplus B,M)$, suppose $\varphi(\tau_1) = \varphi(\tau_2)$. Then $(\tau_1 \circ \iota_A, \tau_1 \circ \iota_B) = (\tau_2 \circ \iota_A, \tau_2 \circ \iota_B)$. Since $\iota_A$ and $\iota_B$ are injective, we have that $\tau_1 = \tau_2$, implying that $\varphi$ is injective. Now let $(f,g) \in \Hom_R(A,M) \times \Hom_R(B,M)$. By the Universal Property of Coproducts, there exists a map $\tau:A \oplus B \rightarrow M$ such that $f = \tau \circ \iota_A$ and $g = \tau \circ \iota_B$. Thus $\varphi(\tau) = (\tau \circ \iota_A, \tau \circ \iota_B) = (f,g)$.

		(b) Since products exist in the category of $R$-modules, let $\pi_A:A \times B \rightarrow A$ and $\pi_B:A \times B \rightarrow B$ be the natural projections. Define $\psi:\Hom_R(A,M \times N) \rightarrow \Hom_R(A,M) \times \Hom_R(A,N)$ by $\psi(f) = (\pi_M \circ f, \pi_N \circ f)$. For $f_1,f_2 \in \Hom_R(A,M\times N)$, we have:
		\begin{align*}
			\psi(f_1+f_2)
			&= \bigl(\pi_M\circ(f_1+f_2),\pi_N\circ(f_1+f_2)\bigr) \\
			&= (\pi_M\circ f_1,\pi_N\circ f_1)
			  +(\pi_M\circ f_2,\pi_N\circ f_2) \\
			&= \psi(f_1)+\psi(f_2).
		\end{align*}
		For $f\in\Hom_R(A,M\times N)$ and $r\in R$, we have:
		\begin{align*}
			\psi(rf)
			&= \bigl(\pi_M\circ(rf),\pi_N\circ(rf)\bigr) \\
			&= r(\pi_M\circ f,\pi_N\circ f) \\
			&= r\psi(f).
		\end{align*}
		Thus $\psi$ is an $R$-module homomorphism. It remains to show that
		$\psi$ is injective and surjective. For $f_1,f_2\in\Hom_R(A,M\times N)$,
		suppose $\psi(f_1)=\psi(f_2)$. Then $\pi_M\circ f_1=\pi_M\circ f_2$ and $\pi_N\circ f_1=\pi_N\circ f_2$.	By the Universal Property of Products, there is a unique map from $A$ to
		$M\times N$ with these two compositions. Hence $f_1=f_2$, implying that
		$\psi$ is injective. Now let
		$(g,h)\in\Hom_R(A,M)\times\Hom_R(A,N)$. By the Universal Property of
		Products, there exists a map $f:A\to M\times N$ such that
		$g=\pi_M\circ f$ and $h=\pi_N\circ f$. Thus $\psi(f)=(\pi_M\circ f,\pi_N\circ f)=(g,h)$. Therefore $\psi$ is surjective, and hence an isomorphism.
	\end{proof}}

\begin{exercise}
Use universal properties to show that there is a canonical morphism
$\theta:\coprod_I M_i\longrightarrow\prod_I M_i$, defined by
$\theta(\sum m_i)=(m_i)$. You should use the fact that in the category of
$R$-modules there is a zero object.
\end{exercise}
	{\color{blue}\begin{proof}
		Since the category of $R$-modules admits a zero object, from the following diagram:
		\begin{center}
		\begin{tikzcd}
		M_i \arrow[r] & 0 \arrow[r] & M_j
		\end{tikzcd}
		\end{center}
		we see that there is a canonical zero morphism between any two modules. For each $j \in I$, define $\delta_j:M_j \rightarrow \prod_{I}M_i$ by:
		\[
			\delta_j(x)=
			\begin{cases}
				x, & i = j \\
				0, & i \neq j
			\end{cases}.
		\] 
		By the Universal Property of Coproducts, there exists a unique map $\theta:\coprod_I M_i \rightarrow \prod_I M_i$ such that the following diagram commutes:
		\begin{center}
\begin{tikzcd}
            & \coprod_I M_i \arrow[ld, "\theta"', dotted]     \\
\prod_I M_i & M_j \arrow[l, "\delta_j"] \arrow[u, "\iota_j"']
\end{tikzcd}		\end{center}
		where $\iota_j$ is the natural inclusion.
	\end{proof}}

\newpage
\begin{exercise}
Let $R$ be a commutative ring with $1$, and suppose we have a short exact
sequence of $R$-modules. A \textit{short exact sequence} of $R$-modules is a diagram of the form
\[
\begin{tikzcd}
1 \arrow[r] & M' \arrow[r, "\iota"] & M \arrow[r, "\pi"] & M'' \arrow[r] & 1
\end{tikzcd}
\]
such that the kernel of each map equals the image of the preceding map, where
applicable. Thus $\iota$ is injective, $\pi$ is surjective, and
$\iota(M')=\ker(\pi)$, so that $M/\iota(M')\simeq M''$. We say the sequence
\textit{splits} if there is a morphism $\sigma:M''\longrightarrow M$ such that
$\pi\circ\sigma=\operatorname{id}_{M''}$. Use the universal property of free
modules to prove that if $M''$ is free then the sequence splits, and then
$M\simeq M'\oplus M''$.
\end{exercise}
	{\color{blue}\begin{proof}
		Since $M''$ is free, let $\{x_i\}_I$ be its basis. Let $e: \{x_i\}_I \rightarrow M''$ be the natural inclusion.	Since $\pi$ is surjective, pick $\{m_i\}_I \subset M$ such that $\pi(m_i) = x_i$ for each $i \in I$. Define $f: \{x_i\}_I \rightarrow M$ by $f(x_i) := \pi^{-1}(e(x_i)) = m_i$. By the Universal Property of Free Products, there exists a unique $R$-linear map $\sigma:M'' \rightarrow M$ such that $f = \sigma \circ e$. Thus:
		\begin{align*}
			\pi(\sigma(x_i))
			&= \pi(\sigma(e(x_i))) \\
			&= \pi(f(x_i)) \\
			&= \pi(m_i) \\
			&= x_i.
		\end{align*}
		Since $\pi \circ \sigma$ is the identity on the basis elements of $M''$, we can conclude that $\pi \circ \sigma = \id_{M''}$.

		Now define $\varphi:M' \oplus M'' \rightarrow M$ by $m' + m'' \mapsto \iota(m') + \sigma(m'')$. Clearly this map is $R$-linear. To show that the map is injective, suppose $\varphi(m' + m'') = 0$. Then $\iota(m') + \sigma(m'') = 0$. Apply $\pi$ to both sides yields $\pi(\iota(m')) + \pi(\sigma(m'')) = 0$. By exactness, $\pi(\iota(m')) = 0$. Since our sequence splits, $\pi(\sigma(m'')) = m''$. Thus $m'' = 0$. Hence $\iota(m') = 0$, and because $\iota$ is injective it must be that $m' = 0$. Thus $m' + m'' = 0$, establishing that $\varphi$ is injective. Now let $m \in M$. Note that $\pi(m - \sigma(\pi(m))) = 0$, hence $m - \sigma(\pi(m)) \in \ker\pi$. Exactness of our sequence gives $\ker\pi = \iota(M')$. So there exists some $m' \in M'$ such that $m - \sigma(\pi(m)) = \iota(m')$. Thus $m = \iota(m') + \sigma(\pi(m)) = \varphi(m',\pi(m))$. Thus $M \cong M' \oplus M''$.

	\end{proof}}

\begin{exercise}
Let $R$ be a commutative ring with $1$. Suppose we have a commutative diagram
of $R$-modules
\[
\begin{tikzcd}
0 \arrow[r] & M' \arrow[r] & M \arrow[r, "f"] \arrow[d, "\wr"] &
M'' \arrow[r] \arrow[d, "\wr"] & 0 \\
0 \arrow[r] & N' \arrow[r] & N \arrow[r, "g"] & N'' \arrow[r] & 0
\end{tikzcd}
\]
Prove that there is an isomorphism $M'\longrightarrow N'$ that makes the
resulting left square commutative.
\end{exercise}


\end{document}
